\answer{ex:20:1}{$t_{\text{각성}}=\tau_r\ln\frac{1-L}{1-H}=16.8$~h, $t_{\text{수면}}=\tau_d\ln\frac HL=5.8$~h, 주기 $22.5$~h; 발진기를 $24$~h로 끌어오는 것은 문턱값의 일주기 흔들림, 즉 위상 고정 루프~\eqref{eq:pll}다.}
\answer{ex:20:2}{(a)~$L=0.111$; $L=0.092$이면 수면은 $9.6$~h 이어졌을 것이다. (b)~$18\%$가 아니라 $22\%$($1/0.82=1.22$).}
\answer{ex:20:3}{$H=\frac{p-1}{p-1/q}=0.623$, $L=H/q=0.093$, 여기서 $p=e^{16/\tau_r}$, $q=e^{8/\tau_d}$. $4$~h 만에 깨우면 위상이 영구히 옮겨 간다(잠드는 시각이 $7.2$~h 빨라진다). 문턱값이 오르내리면 위상은 2--3일 안에 돌아온다.}
\answer{ex:20:4}{(a)~$r_\pm=\frac12\bigl(1\pm\sqrt{1-4k/w}\bigr)=0.947,\ 0.053$; $a_{\text{상}}=3.51$, $a_{\text{하}}=0.49$. (b)~활동 $\approx0.33$~s, 침묵 $\approx0.59$~s, 주기 $\approx0.93$~s, 활동 비율 $0.36$.}
\answer{ex:20:5}{$a_{\text{상}}/r_+<b<a_{\text{하}}/r_-$: $3.71<b<9.20$. \nt{ACET}이 $b$를 $3.71$ 아래로 낮추면 교점이 위 가지로 옮겨 가고, 활동 상태($r\approx1$)가 안정해진다.}
\answer{ex:20:6}{$D=24$, $32$, $40$, $48$, $64$에서 $4.9$, $6.8$, $8.8$, $5.5$, $8.9$~h: 길이를 정하는 것은 빚이 아니라 잠드는 시각의 일주기 위상이다.}
\answer{ex:20:7}{B 뒤 A에 대한 오차는 평균 $0.51$($20$개 세트에서 $0.19$부터 $1.08$까지; 응답이 0이면 $1$). 섞어서 배우면 두 오차 모두 $10^{-4}$보다 작다.}
\answer{ex:20:8}{열린 문제. 균일한 스케일링은 가중치의 비를 보존하지만, 문턱 뉴런의 응답은 문턱값도 같은 비율로 스케일링될 때만 보존된다. 섞은 재생은 비를 바꾼다. 큰 접촉이 바뀌지 않는다는 사실은 순수하게 균일한 스케일링과 모순된다.}
